\documentclass{article}
\usepackage[T1]{fontenc}
\usepackage{lmodern}
\usepackage{graphicx}
\usepackage{amsmath,amssymb}
\usepackage[a4paper,margin=2cm]{geometry}
\usepackage[absolute,overlay]{textpos}
\setlength{\TPHorizModule}{1mm}
\setlength{\TPVertModule}{1mm}
\usepackage[indonesian]{babel}
\usepackage{hyperref}
\setlength{\emergencystretch}{2em}
% unit: Halaman 14

\begin{document}

% === Halaman 14 (llm) ===

\section*{Contoh pengiriman pesan 2:}
\begin{itemize}
    \item Misalkan Bob akan mengirim plainteks $M = \text{'HELLO ALICE'}$ kepada Alice
    \item Dengan memisalkan $\text{A} = 00, \text{B} = 01, \dots, \text{Z} = 25$, maka plainteks dikodekan ke dalam \textit{integer} (spasi diabaikan) menjadi
    \[ m = 070411111140011080204 \]
\end{itemize}

Nilai $m$ ini sangat besar untuk dipangkatkan dan tidak berada di dalam selang $[0, 3337-1]$, oleh karena itu pecah $m$ menjadi blok-blok pesan yang lebih kecil, misalnya blok 4 digit:

\begin{align*}
    m_1 &= 0704 & m_4 &= 1108 \\
    m_2 &= 1111 & m_5 &= 0204 \\
    m_3 &= 1400 & \\
\end{align*}

(Perhatikan, $m_i$ masih valid karena terletak di dalam selang $[0, 3337 - 1])$

\end{document}
